%! Inverse Functions — Unit 1, Lesson 1.6 — Group Activity (ANSWER KEY)
\documentclass[10pt]{article}
\usepackage{precalculus-article}
\usepackage{precalculus-key}
\newcommand{\TallMath}[1]{$\displaystyle #1\rule[-1.4em]{0pt}{3.2em}$}

\begin{document}

\pageheader{Unit 1, Lesson 1.6}{Group Activity --- Answer Key}

\vspace{0.12in}

\begin{scenariobox}[The Charging Station]{plumlight}
\small
An electric car pulls into a fast charger with 10 miles of range left. After
$t$ hours on the charger its range is

\begin{center}
$R(t) = 25t + 10$ \quad miles.
\end{center}

The dashboard reports one thing only: \textbf{miles of range}. But the driver
standing next to the car is always asking the other question ---
\textbf{how long do I have to sit here?} Every question in this activity is
that question.

\begin{center}
\begin{tikzpicture}[x=1.0cm, y=0.30cm]
  \draw[linegray, very thin, xstep=1, ystep=2] (0,0) grid (4,12);
  \draw[->, thick] (0,0) -- (4.7,0) node[below right] {\small $t$ (h)};
  \draw[->, thick] (0,0) -- (0,13.6) node[above] {\small $R(t)$ (mi)};
  \foreach \x in {1,2,3,4} \node[below, font=\small] at (\x,0) {\x};
  \foreach \y/\lab in {2/20, 4/40, 6/60, 8/80, 10/100, 12/120}
    \node[left, font=\small] at (0,\y) {\lab};
  \draw[very thick, plum] (0,1) -- (4,11);
  \fill[plum] (0,1) circle (2.5pt);
  \fill[plum] (2,6) circle (2.5pt);
  \fill[plum] (4,11) circle (2.5pt);
\end{tikzpicture}
\end{center}

\renewcommand{\arraystretch}{1.3}
\begin{center}
\begin{tabular}{|c|c|c|c|c|c|}
\hline
$t$ (hours) & 0 & 1 & 2 & 3 & 4 \\
\hline
$R(t)$ (miles) & 10 & 35 & 60 & 85 & 110 \\
\hline
\end{tabular}
\end{center}
\end{scenariobox}

\vspace{0.1in}

\boxguard[26]
\begin{tcolorbox}[colback=white, colframe=black!40,
    title=\textbf{Tier R --- Remediate},
    fonttitle=\bfseries, arc=2mm, left=3mm, right=3mm, top=2mm, bottom=2mm]

\begin{enumerate}[itemsep=6pt]
  \item Read the scenario table forwards. After 3 hours the range is
        \ans{85} miles.

  \item Now read it \textbf{backwards} --- the way the driver reads it. The
        dashboard says 85 miles. How long has the car been charging?

        \ans{3} hours

  \item Build $R^{-1}$ by \textbf{trading the two rows} of the scenario table.
        Nothing is computed --- the answers become the questions.

        \smallskip
        \renewcommand{\arraystretch}{1.4}
        \begin{center}
        \begin{tabular}{|c|c|c|c|c|c|}
        \hline
        $x$ (miles) & 10 & 35 & 60 & 85 & 110 \\
        \hline
        $R^{-1}(x)$ (hours) & \ans{0} & \ans{1} & \ans{2} & \ans{3} & \ans{4} \\
        \hline
        \end{tabular}
        \end{center}

  \item Fill in the units, and watch them trade places.

        \smallskip
        \renewcommand{\arraystretch}{1.4}
        \begin{center}
        \begin{tabularx}{0.9\linewidth}{@{} >{\bfseries}p{0.16\linewidth} X X @{}}
        \hline
                 & \textbf{takes in} & \textbf{gives back} \\
        \hline
        $R$      & \ans{hours} & \ans{miles} \\
        \hline
        $R^{-1}$ & \ans{miles} & \ans{hours} \\
        \hline
        \end{tabularx}
        \end{center}

  \item In one sentence, say what $R^{-1}$ does for the driver.

        \vspace{0.05in}
        \ansline{It reads the number on the dashboard and gives back the question --- how long the car has been on the charger.}
\end{enumerate}
\end{tcolorbox}

\vspace{0.1in}

\boxguard[26]
\begin{tcolorbox}[colback=white, colframe=black!40,
    title=\textbf{Tier A --- Approaching Proficiency},
    fonttitle=\bfseries, arc=2mm, left=3mm, right=3mm, top=2mm, bottom=2mm]

The table only answers the five questions printed on it. A formula answers all
of them.

\begin{enumerate}[itemsep=8pt]
  \item Find $R^{-1}(x)$ by swapping and solving.

        \begin{work}
                        y &= 25x + 10 \\
                        x &= 25y + 10 \\
                   x - 10 &= 25y \\
          \frac{x-10}{25} &= y
        \end{work}

        $R^{-1}(x) = $ \ans{$\dfrac{x-10}{25}$}

  \item Check it against the table: $R^{-1}(85) = $ \ans{3} hours.

  \item The driver needs 160 miles of range before leaving.

        \begin{enumerate}[label=(\alph*), itemsep=4pt]
          \item $R^{-1}(160) = $ \ans{6} hours.
          \item Say in one sentence why $R^{-1}$ --- and not $R$ --- is the
                function that answers this.

                \vspace{0.05in}
                \ansline{The driver was handed an answer (160 miles) and wants the question back; $R$ only runs hours forward into miles.}
        \end{enumerate}

  \item Verify the pair in \textbf{both} directions.

        \begin{work}
          R\big(R^{-1}(x)\big) &= 25\left(\frac{x-10}{25}\right) + 10 \\
                               &= (x - 10) + 10 \\
                               &= x \\
          R^{-1}\big(R(x)\big) &= \frac{(25x + 10) - 10}{25} \\
                               &= \frac{25x}{25} \\
                               &= x
        \end{work}

        Both came out to \ans{$x$}, which is the \ans{identity} function.

  \item \textbf{The trap, in numbers.} The dashboard reads 60 miles.

        \begin{enumerate}[label=(\alph*), itemsep=4pt]
          \item $R^{-1}(60) = $ \ans{2} hours.
          \item $R(60) = $ \ans{1510}, so $\dfrac{1}{R(60)} = $
                \ans{$1/1510$}.
          \item One of those two numbers is a charging time. Which one, and
                what is the other one even measuring?

                \vspace{0.05in}
                \ansline{$R^{-1}(60) = 2$ hours is the time; $1/1510$ came from charging for 60 \textit{hours} and flipping the miles --- it measures nothing.}
        \end{enumerate}
\end{enumerate}
\end{tcolorbox}

\vspace{0.1in}

\boxguard[26]
\begin{tcolorbox}[colback=white, colframe=black!40,
    title=\textbf{Tier E --- Extension},
    fonttitle=\bfseries, arc=2mm, left=3mm, right=3mm, top=2mm, bottom=2mm]

\begin{enumerate}[itemsep=8pt]
  \item \textbf{The battery warms up, then cools.} During a 6-hour charge the
        battery temperature is $T(t) = (t-3)^2 + 50$ degrees F.

        \begin{work}
          T(1) &= (1-3)^2 + 50 \\
               &= 54 \\
          T(5) &= (5-3)^2 + 50 \\
               &= 54
        \end{work}

        \begin{enumerate}[label=(\alph*), itemsep=4pt]
          \item $T(1) = $ \ans{54} and $T(5) = $ \ans{54}, so the
                answer $54^\circ$ came from \ans{two} different questions.
          \item Therefore $T$ is not \ans{one-to-one} on $0 \le t \le 6$, and a
                thermometer reading alone \ans{cannot} tell you the time.
          \item \textbf{Repair it.} Keep only $t \ge 3$ and invert on that
                piece.

                \begin{work}
                      y &= (t-3)^2 + 50 \\
                 y - 50 &= (t-3)^2 \\
              \sqrt{y-50} &= t - 3 \\
              3 + \sqrt{y-50} &= t
                \end{work}

                $T^{-1}(y) = $ \ans{$3 + \sqrt{y-50}$}

          \item Test it: $T^{-1}(54) = $ \ans{5}. Which of the two times
                from part (a) did the restriction keep? \ans{$t = 5$}
        \end{enumerate}

  \item \textbf{Dollars from miles.} The station bills $C(t) = 8t + 2$ dollars
        for $t$ hours on the charger. The driver wants cost as a function of
        the number on the \textit{dashboard} --- and no meter reports that.

        \begin{work}
          C\big(R^{-1}(x)\big) &= 8\left(\frac{x-10}{25}\right) + 2 \\
                               &= 0.32(x - 10) + 2 \\
                               &= 0.32x - 1.2
        \end{work}

        \begin{enumerate}[label=(\alph*), itemsep=4pt]
          \item The chain runs $R^{-1}$ first, then $C$. The handoff between
                them is a number of \ans{hours}.
          \item $C\big(R^{-1}(x)\big) = $ \ans{$0.32x - 1.2$}
          \item Check at 60 miles: the composite gives \$\ans{18}, and
                the table says 60 miles is 2 hours, so $C(2) = $
                \$\ans{18}.
          \item Why can the chain not be run the other way, as
                $R^{-1}\big(C(t)\big)$?

                \vspace{0.05in}
                \ansline{$C$ hands over dollars, and $R^{-1}$ only accepts miles --- the handoff has the wrong units, so there is nothing to compute.}
        \end{enumerate}
\end{enumerate}
\end{tcolorbox}

\end{document}
